% PARKED CONTENT -- not \input by main.tex.
% Cut from the results section when it was cleared; kept here because Sections_zh
% is not under version control and this text is not recoverable from git.

% ---------------------------------------------------------------- per-task panels
\begin{figure}[t]
  \centering
  \includegraphics[width=\linewidth]{fig_per_task}
  \caption{\textbf{Every configuration, every task, against that task's baseline.}
  One panel per task; each point is one of the 29 (model, effort) configurations,
  coloured by model, and the dashed line is the task's baseline. The ten tasks
  share neither scale nor direction --- two are minimised, the metrics run from a
  perplexity in the tens to a pass rate in $[0,1]$ --- so each panel carries its
  own axis, and what is comparable across panels is only which side of the line a
  configuration lands on. Missing points are configurations that returned no
  trainable model; they count against the system rather than being dropped.}
  \label{fig:per_task}
\end{figure}

% ---------------------------------------------------------------- patch families
\subsection{What a submission changes}
\label{sec:families}

A score says whether a submission worked. It does not say whether the agent designed an algorithm or
tuned one, and those are the two outcomes this paper needs to tell apart. Nor does the size of the
change settle it: a one-line diff can replace a learning rule, and a thousand-line refactor can leave
the training procedure exactly as it was. So the distinction has to be read off the submitted code
rather than off the score, and we read it by classifying every submission according to which part of
the training procedure it modifies.

Eight families cover what the submissions do, and they fall on two sides of that line. Four change
\emph{how the run is executed}: budget, checkpointing and logging; optimisation knobs; selection and
averaging over checkpoints; and capacity or adapter placement. Four change \emph{how the model
learns}: surgery on the objective; added supervision; substitution of the learning rule; and
intervention in the data the procedure trains on. The families are not exclusive --- a submission
matches three of them on average, since changing an objective usually drags a knob along with it ---
so we report how many submissions reach a side rather than partitioning them into one family each.

Of the 280 submissions, 17 changed nothing that could be classified. Among the 263 that did, 122
reach the learning side and 141 stay entirely on the run side. That split is the instrument behind
\S\ref{sec:results}, and it is worth being clear about what it can and cannot support: it records
which part of the procedure a submission touched, not whether the touch was a good idea. A submission
that rewrites an objective badly and one that rewrites it well are both on the learning side. That is
the point --- it lets us ask, separately from whether an agent won, whether it was even working on
the part of the system where the compounding lives.

% ---------------------------------------------------------------- cut from the analysis section
\subsection{Agents optimize the scoring, not the intent}
\label{sec:hacking}

None of the ten metrics is ours: freezing a real repository means inheriting the objective it was
already optimising. An objective that behaves well over the range its own authors worked in need not
survive an agent pushing deliberately on it, and two places in this suite show what that looks
like.

\textbf{Agents bypass reinforcement learning and learn from the solver.} The agentic-RL repository
ships an optimal solver inside its environment tooling. It sits on no path the evaluator uses, but it
is readable, and four of the five configurations that scored perfectly used it to generate
supervision; 11 of the 29 submissions call it directly. This breaks no rule: what the task forbids is
importing external boards, demonstrations, trajectories or weights, and training on the final seeds,
and the solver is inside the training container. The fifth perfect configuration is the control ---
it opened the same file, did not read as far as the section holding that function, wrote its own
breadth-first solver instead, and reached the same perfect score. A perfect score in this column
therefore certifies that a system recognised the task as one for imitation learning, not that it
trained reinforcement learning better, which also explains the column's shape: five cells at the
ceiling and a median still at $0.144$.

\textbf{Agents hedge with the three checkpoint slots.} A cell is scored on the best of three retained
checkpoints, and three slots can be used to hedge. On one-shot pruning, one submission gave its three
slots to three finished products rather than three moments: pure pruning with no training at all,
perplexity $21.8$; a fine-tune stopped halfway, $14.6$; and the full decay schedule, $14.3$. It
completed and published the first in about ten minutes, spent the remaining time on the other two,
and left retries in the script. This is a fully compliant and rather sophisticated submission ---
bank a deliverable that is certain to exist, then gamble the rest of the budget on something better.
It also means the three checkpoints of a configuration need not be three points on one training
curve: best-of-three permits choosing among methods as much as choosing along a curve.

What both have in common is that the agent treats the task as an object to be optimized rather
than as an intent to be served. That is the capability we want it to have; the same capability finds
the seams in whatever is being scored. Validity checks therefore have to be written per task, and
they have to block rather than advise.
